\section{Os Axiomas}
\index{axioma}
\index{ZF}
\index{ZFC}
\index{Z}
\index{ZC}

A teoria dos conjuntos moderna é usualmente desenvolvida em uma linguagem cujo único símbolo não lógico primitivo é $\in$.
Os princípios não lógicos que regem essa linguagem são os axiomas.
O sistema hoje chamado $ZFC$ tem sua origem na axiomatização de Zermelo de 1908, concebida para evitar os paradoxos da compreensão irrestrita \cite{zermelo1908untersuchungen}.
Em 1922, Fraenkel e, independentemente, Skolem observaram que a teoria de Zermelo era insuficiente para justificar certas construções recursivas elementares e propuseram o princípio que deu origem ao atual Esquema da Substituição \cite{fraenkel1922grundlagen,skolem1923bemerkungen}.
Pouco depois, von Neumann introduziu o Axioma da Regularidade, também chamado de Axioma da Fundação \cite{vonneumann1925axiomatisierung}.
A reformulação de Zermelo de 1930 consolidou uma apresentação já bastante próxima das axiomatizações contemporâneas \cite{zermelo1930grenzzahlen}.

Neste texto, enunciaremos desde já todos os axiomas de $ZFC$.
Não apresentaremos nenhuma explicação sobre eles nessa seção.
Por ora, queremos apenas introduzir a terminologia que usaremos para se referir aos axiomas.

Os axiomas serão introduzidos e estudados em detalhe ao longo dos capítulos iniciais do texto.
Em particular, o Axioma do Infinito será estudado apenas no Capítulo~\ref{chap:numerosNaturais}, o Esquema da Substituição no Capítulo~\ref{chap:numerosOrdinais} e o Axioma da Escolha no Capítulo~\ref{chap:axiomaDaEscolha}.

\subsection{Enunciado dos axiomas}

A seguir, enunciaremos os axiomas de ZFC na forma em que eles serão usados no texto.
Nos enunciados abaixo, usaremos as abreviações lógicas e notacionais já introduzidas anteriormente,
como quantificadores restritos, $\neq$ e $\exists!$.
Se $s, t$ são termos da linguagem, $s \notin t$ é uma abreviação de $\neg (s \in t)$.

Existem, na literatura, diversas reformulações equivalentes desta lista de axiomas.
Optamos por esta lista em particular pela simplicidade de sua apresentação e por sua adequação ao desenvolvimento que faremos ao longo do texto.

\begin{displayaxiom}{Axioma da Extensão}\index{axioma!da extensão}
    \[
        \forall x \forall y\, \Bigl(\bigl(\forall z\, (z \in x \leftrightarrow z \in y)\bigr) \rightarrow x = y\Bigr).
    \]
\end{displayaxiom}

\begin{displayaxiom}{Esquema de Axiomas da Compreensão}\index{axioma!esquema da compreensão}\index{axioma!esquema da separação|see{esquema da compreensão}}
    Para cada fórmula $\phi(x, A, t_1, \ldots, t_n)$ sem $y$ livre, a seguinte sentença é um axioma:
    \[
        \forall t_1 \ldots \forall t_n \forall A \exists y \forall x\, \Bigl(x \in y \leftrightarrow (x \in A \land \phi(x, A, t_1, \ldots, t_n))\Bigr).
    \]
\end{displayaxiom}

\begin{displayaxiom}{Axioma do Par}\index{axioma!do par}
    \[
        \forall x \forall y \exists z\, (x \in z \land y \in z).
    \]
\end{displayaxiom}

\begin{displayaxiom}{Axioma da União}\index{axioma!da união}
    \[
        \forall x \exists y \forall z\, \bigl(\exists w\, (w \in x \land z \in w) \rightarrow z \in y\bigr).
    \]
\end{displayaxiom}

\begin{displayaxiom}{Axioma da Potência}\index{axioma!da potência}
    \[
        \forall x \exists y \forall z\, \Bigl((\forall w\, (w \in z \rightarrow w \in x)) \rightarrow z \in y\Bigr).
    \]
\end{displayaxiom}

\begin{displayaxiom}{Axioma do Infinito}\index{axioma!do infinito}
    \[
        \begin{aligned}
            \exists I\, \Bigl(
                &\exists z\, \bigl(\forall w\, (w \notin z) \land z \in I\bigr) \\
                &\land \forall x\, \bigl(
                    x \in I
                    \rightarrow
                    \exists y\, \bigl(
                        y \in I
                        \land
                        \forall t\, (t \in y \leftrightarrow (t \in x \lor t = x))
                    \bigr)
                \bigr)
            \Bigr).
        \end{aligned}
    \]
\end{displayaxiom}

\begin{displayaxiom}{Esquema de Axiomas da Substituição}
    \index{axioma!esquema da substituição}
    Para cada fórmula $\phi(x, y, A, t_1, \ldots, t_n)$ sem $B$ livre, a seguinte sentença é um axioma:
    \[
        \begin{aligned}
            \forall t_1 \ldots \forall t_n \forall A\, \Bigl(
                &\forall x \in A\, \exists! y\, \phi(x, y, A, t_1, \ldots, t_n) \\
                &\rightarrow \exists B\, \forall y\, \bigl(\exists x \in A\, \phi(x, y, A, t_1, \ldots, t_n) \rightarrow y \in B\bigr)
            \Bigr).
        \end{aligned}
    \]
\end{displayaxiom}

\begin{displayaxiom}{Axioma da Regularidade (ou da Fundação)}\index{axioma!da regularidade}\index{axioma!da fundação|see{da regularidade}}
    \[
        \forall x\, \Bigl((\exists y\, (y \in x)) \rightarrow \exists y\, \bigl(y \in x \land \neg \exists z\, (z \in x \land z \in y)\bigr)\Bigr).
    \]
\end{displayaxiom}

\begin{displayaxiom}{Axioma da Escolha}\index{axioma!da escolha}
    \[
        \begin{aligned}
            \forall A\, \Bigl(
                &\Bigl(
                    \forall x \in A\, \exists y\, (y \in x) \\
                    &\land \forall x \in A\, \forall y \in A\, \bigl(x \neq y \rightarrow \neg \exists z\, (z \in x \land z \in y)\bigr)
                \Bigr) \\
                &\rightarrow \exists C\, \forall x \in A\, \exists! y\, (y \in x \land y \in C)
            \Bigr).
        \end{aligned}
    \]
\end{displayaxiom}

No Capítulo~\ref{chap:axiomaDaEscolha}, estudaremos esse axioma e algumas de suas formulações equivalentes.

\subsection{Nomenclatura}
    As letras $Z$ e $F$ em $ZF$ e $ZFC$ remetem a Zermelo e Fraenkel.
    Já $C$ remete ao Axioma da Escolha.
    Abaixo, $P$ remete ao Axioma da Potência.
\begin{metadefinition}[Nomes das teorias]
    Ao longo do livro, adotaremos as seguintes nomenclaturas.
    \label{mdef:nomesDasTeorias}
    \begin{itemize}
        \item $Z$ é a teoria formada pelos seguintes axiomas e esquemas: extensão, Compreensão, Par, União, Potência, Infinito e Regularidade.
        \item $ZF$ é a teoria $Z$ acrescida do Esquema da Substituição.
        \item $ZC$ é a teoria $Z$ acrescida do Axioma da Escolha.
        \item $ZFC$ é a teoria $ZF$ acrescida do Axioma da Escolha.
        \item Se $T$ é uma das teorias acima, $T^-$ é a teoria obtida ao se remover de $T$ o Axioma da regularidade.
        \item Se $T$ é uma das teorias acima, $T-P$ é a teoria obtida ao se remover de $T$ o Axioma da Potência.
        \item $BST$, ou \emph{Basic Set Theory}, é a teoria formada pelos seguintes axiomas e esquemas: Extensão, Compreensão, União, Par, e Potência fraca.
    \end{itemize}

    O Axioma da Potência Fraca não é um axioma padrão da literatura.
    Neste texto, ele será utilizado para fins didáticos e organizacionais.
    Na presença dos demais axiomas considerados nesta seção, ele é uma versão muito mais fraca do
    Axioma da Potência e também decorre do Esquema da Substituição.
\end{metadefinition}

\begin{displayaxiom}{Axioma da Potência Fraca}\index{axioma!da potência fraca}
        \[
            \begin{aligned}
                \forall x\, \exists y\, \forall a\, \forall b\, \Bigl(
                    &(a \in x \land b \in x) \\
                    &\rightarrow \exists z\, \bigl(
                        z \in y
                        \land
                        \forall w\, (w \in z \leftrightarrow (w = a \lor w = b))
                    \bigr)
                \Bigr).
            \end{aligned}
        \]
    \end{displayaxiom}
